%
% channel-check-ackopen.tex  (FROZEN NEGATIVE CONTROL: do not edit)
%
% Acknowledge leaves the token pending: a repeated channel_ack is relayed again.
%
% Z Specification for the channel check: the request biff doctor sends to a
% session, the one channel event it makes, and the acknowledgement the model
% returns through the channel_ack tool (design-channels.md section 5, DES-087).
%

\documentclass[a4paper,10pt,fleqn]{article}
\usepackage[margin=1in]{geometry}
\usepackage{fuzz}

\begin{document}

\title{Biff Channel Check: A Z Specification}
\author{Request, Token, Event and Acknowledgement}
\date{November 2026}
\maketitle

\section{Introduction}

\texttt{biff doctor} proves, one session at a time, that channel events reach a
model.  The session's server accepts a request that names a one-time reply
subject, makes a fresh token, sends one channel event carrying it, and accepts
the token back once through \texttt{channel\_ack}.  The server accepts at most one
request a minute.  The check only reports: no operation below touches what is
delivered, and the specification has no variable for delivery.

\section{Basic Types}

A \emph{token} is the 128-bit hex string of one check.  A \emph{subject} is a
NATS subject.  The reply subject of a request is a fresh inbox name; a subject
that is not an inbox is \emph{foreign}.

\begin{zed}
  [TOKEN, SUBJECT]
\end{zed}

\begin{zed}
  Report ::= live | notAccepted
\end{zed}

\begin{axdef}
  checkSeconds : \nat_1 \\
  tickSeconds : \nat_1 \\
  maxClock : \nat \\
  inboxes : \power SUBJECT
\where
  checkSeconds = 60 \\
  tickSeconds = 30 \\
  maxClock = 120
\end{axdef}

The clock moves in steps of \emph{tickSeconds}, half a window, and stops at
\emph{maxClock}: the bound that lets ProB exhaust the state space.

\section{State}

Time is a natural number of seconds.  At most one token is pending: a new
accepted request retires the older one.  A token that was acknowledged is
\emph{spent} and never matches again.  The times of accepted requests are
recorded to state the once-a-minute bound as an invariant.

\begin{schema}{CheckState}
  clock : \nat \\
  acceptedAt : \power \nat \\
  pending : \power TOKEN \\
  unsent : \power TOKEN \\
  sent : \power TOKEN \\
  spent : \power TOKEN \\
  replyOf : TOKEN \pfun SUBJECT \\
  delivered : TOKEN \pfun (0 \upto 2) \\
  emitted : TOKEN \pfun (0 \upto 2) \\
  relayed : TOKEN \pfun (0 \upto 2)
\where
  clock \leq maxClock \\
  \# pending \leq 1 \\
  unsent \subseteq pending \\
  pending \cap spent = \emptyset \\
  unsent \cap sent = \emptyset \\
  pending \subseteq \dom replyOf \\
  \ran replyOf \subseteq inboxes \\
  \forall a, b : acceptedAt | a \neq b @ a + checkSeconds \leq b \lor b +
    checkSeconds \leq a
\end{schema}

The predicates of $CheckState$ are, to ProZ, guards of every operation: a step
that would break one is not enabled, and the model checker cannot report it.  The
properties of the daemon hop are therefore stated in a schema of their own, which
$Invariant$ installs as the checked state invariant, so that a step breaking one
is enabled, reaches the bad state, and is reported.

\begin{schema}{HopInv}
  CheckState
\where
  \dom emitted = \dom delivered \\
  \dom relayed = \dom delivered \\
  \forall t : \dom emitted @ emitted~t \leq 1 \\
  \forall t : \dom relayed @ relayed~t \leq 1 \\
  \forall t : sent @ emitted~t = 1
\end{schema}

The token is single-use across the daemon hop: a request is turned into at most
one channel event whatever the number of deliveries ($emitted~t \leq 1$), the
acknowledgement is relayed to the doctor at most once ($relayed~t \leq 1$), and a
token that was sent was emitted exactly once.

The replay case is a reachability goal: $spent \neq \emptyset \land \exists t
@ t \in spent \land delivered~t = 2 \land emitted~t = 1$, a request frame delivered
again after its token was acknowledged, with still exactly one event made.

\begin{schema}{Invariant}
  HopInv
\end{schema}

\begin{schema}{Init}
  CheckState~'
\where
  clock' = 0 \\
  acceptedAt' = \emptyset \\
  pending' = \emptyset \\
  unsent' = \emptyset \\
  sent' = \emptyset \\
  spent' = \emptyset \\
  replyOf' = \emptyset \\
  delivered' = \emptyset \\
  emitted' = \emptyset \\
  relayed' = \emptyset
\end{schema}

\section{Operations}

\subsection{Time}

\begin{schema}{Tick}
  \Delta CheckState
\where
  clock < maxClock \\
  clock' = clock + tickSeconds \\
  acceptedAt' = acceptedAt \\
  pending' = pending \\
  unsent' = unsent \\
  sent' = sent \\
  spent' = spent \\
  replyOf' = replyOf \\
  delivered' = delivered \\
  emitted' = emitted \\
  relayed' = relayed
\end{schema}

A session that is asked nothing is not stuck: it stays as it is.

\begin{schema}{Idle}
  \Xi CheckState
\end{schema}

\subsection{Request}

A request is accepted when its reply is an inbox and no request was accepted in
the last \emph{checkSeconds}.  Its token is fresh: never pending, sent or spent
before.  It retires the older pending token.

\begin{schema}{Request}
  \Delta CheckState \\
  reply? : SUBJECT \\
  token? : TOKEN
\where
  reply? \in inboxes \\
  \forall a : acceptedAt @ a + checkSeconds \leq clock \\
  token? \notin pending \cup sent \cup spent \cup \dom replyOf \\
  clock' = clock \\
  acceptedAt' = acceptedAt \cup \{ clock \} \\
  pending' = \{ token? \} \\
  unsent' = \{ token? \} \\
  sent' = sent \\
  spent' = spent \\
  replyOf' = \{ token? \mapsto reply? \} \\
  delivered' = delivered \oplus \{ token? \mapsto 0 \} \\
  emitted' = emitted \oplus \{ token? \mapsto 0 \} \\
  relayed' = relayed \oplus \{ token? \mapsto 0 \}
\end{schema}

A request that names a foreign subject, or arrives inside the window, changes
nothing and does not extend the window.

\begin{schema}{RejectRequest}
  \Xi CheckState \\
  reply? : SUBJECT
\where
  reply? \notin inboxes \lor
    (\exists a : acceptedAt @ clock < a + checkSeconds)
\end{schema}

\subsection{The daemon hop}

With biffd the request does not reach the session's server directly.  The
daemon publishes it, carrying the token, to the session's check subject, and the
session's client receives it as an event frame.  Delivery is at least once: a
reconnect redelivers a frame the client may already have seen.  $delivered~t$
counts how many times the request for $t$ reached the client, and at most two are
modelled.

\begin{schema}{DeliverRequest}
  \Delta CheckState \\
  token? : TOKEN
\where
  token? \in \dom delivered \\
  delivered~token? < 2 \\
  clock' = clock \\
  acceptedAt' = acceptedAt \\
  pending' = pending \\
  unsent' = unsent \\
  sent' = sent \\
  spent' = spent \\
  replyOf' = replyOf \\
  delivered' = delivered \oplus \{ token? \mapsto delivered~token? + 1 \} \\
  emitted' = emitted \\
  relayed' = relayed
\end{schema}

The replay does not depend on the token being pending: the transport may redeliver
a request frame after the acknowledgement closed the token, or after a newer
request retired it, so $DeliverRequest$ is enabled for any token that was ever
requested, up to the two deliveries modelled.

A request that reaches the client again does not make a second event: $SendEvent$
needs the token in $unsent$ and takes it out, so $emitted~t \leq 1$ however
often the frame is delivered.  The acknowledgement the model returns
(\texttt{channel\_ack}) travels the other way, through the client to the daemon;
the daemon relays it once, to the doctor's reply subject $replyOf~t$, and the
token is spent: $relayed~t \leq 1$ however often the acknowledgement frame
arrives.  The three counters, $delivered$, $emitted$ and $relayed$, are ghosts.

\subsection{The one event}

\begin{schema}{SendEvent}
  \Delta CheckState \\
  token? : TOKEN
\where
  token? \in unsent \\
  delivered~token? \geq 1 \\
  clock' = clock \\
  acceptedAt' = acceptedAt \\
  pending' = pending \\
  unsent' = unsent \setminus \{ token? \} \\
  sent' = sent \cup \{ token? \} \\
  spent' = spent \\
  replyOf' = replyOf \\
  delivered' = delivered \\
  emitted' = emitted \oplus \{ token? \mapsto emitted~token? + 1 \} \\
  relayed' = relayed
\end{schema}

\subsection{Acknowledgement}

\begin{schema}{Acknowledge}
  \Delta CheckState \\
  token? : TOKEN \\
  reply! : SUBJECT \\
  report! : Report
\where
  token? \in pending \\
  reply! = replyOf~token? \\
  report! = live \\
  clock' = clock \\
  acceptedAt' = acceptedAt \\
  pending' = pending \\
  unsent' = unsent \\
  sent' = sent \\
  spent' = spent \\
  replyOf' = replyOf \\
  delivered' = delivered \\
  emitted' = emitted \\
  relayed' = relayed \oplus \{ token? \mapsto relayed~token? + 1 \}
\end{schema}

A wrong or repeated token changes nothing and says so.

\begin{schema}{RejectAck}
  \Xi CheckState \\
  token? : TOKEN \\
  report! : Report
\where
  token? \notin pending \\
  report! = notAccepted
\end{schema}

\section{What the implementation owes}

\begin{description}
\item[C1] \texttt{ChannelCheck.request} is \emph{Request} or
  \emph{RejectRequest}; the clock is injected, so a test moves time without
  sleeping.
\item[C2] \texttt{ChannelEmitter} sends the event of \emph{SendEvent} once per
  accepted request, and only with the client initialised.
\item[C3] \texttt{ChannelCheck.acknowledge} is \emph{Acknowledge} or
  \emph{RejectAck}, comparing in constant time.
\item[C4] Nothing here writes to, or reads from, the arrival records or the
  rate window: the check never changes delivery.
\end{description}

\end{document}
